\chapter{اللقاء العشرون: تابع تفسير سورة البقرة (من قوله تعالى: \textit{يَا أَيُّهَا الَّذِينَ آمَنُوا اسْتَعِينُوا بِالصَّبْرِ وَالصَّلَاةِ})}

\section{مقدمة: استئناف التفسير من آيات الصبر والابتلاء}
السلام عليكم ورحمة الله وبركاته. الحمد لله رب العالمين، وأشهد أن لا إله إلا الله وحده لا شريك له، وأشهد أن محمداً عبده ورسوله. اللهم صل على محمد وعلى آل محمد كما صليت على آل إبراهيم إنك حميد مجيد، اللهم بارك على محمد وعلى آل محمد كما باركت على آل إبراهيم إنك حميد مجيد.

استأنف الشيخ هذا اللقاء من قوله تعالى: \begin{ayah}يَا أَيُّهَا الَّذِينَ آمَنُوا اسْتَعِينُوا بِالصَّبْرِ وَالصَّلَاةِ إِنَّ اللَّهَ مَعَ الصَّابِرِينَ\end{ayah}، ونبه من أول الدرس إلى أمر منهجي مهم في صنيع \rawi{الطبري}: وهو أنه يعتني بسياق الآيات وسبب نزولها، لكنه لا يحبس دلالة الآية في ذلك السبب وحده، بل يجمع بين معرفة ما نزلت فيه الآية أولاً، ثم ما تتنزل عليه من المعاني والأحكام.

\begin{fawaid}[الجمع بين التنزيل والتنزل]
من الخطأ أن تُقطع الآية عن سبب نزولها وسياقها، كما أنه من الخطأ أن تُحبس دلالتها في ذلك السبب وحده. والمنهج المحكم أن نبدأ بـ ``التنزيل'' أي: ما نزلت فيه الآية أولاً، ثم ننظر في ``التنزل'' أي: ما يدخل في معناها بعد ذلك. وهذا من المعاني التي يكثر ملاحظتها في تفسير \rawi{الطبري}.
\end{fawaid}

\begin{fawaid}[وحدة السياق من مفاتيح الفهم]
من الدقائق التي نبه عليها الشيخ هنا أن كثيراً من الناس يقرأ الآية مفصولة عما قبلها، مع أن الطبري يلحظ اتصال المقاطع القرآنية بعضها ببعض؛ فالأمر بالصبر والصلاة هنا متصل بفتنة تحويل القبلة وما تبعها من اعتراض السفهاء، ثم هو مع ذلك باق على عمومه في سائر أبواب الطاعة والابتلاء.
\end{fawaid}

\separator

\section{تأويل (استعينوا بالصبر والصلاة إن الله مع الصابرين)}
\isnad{قال أبو جعفر:} إن هذه الآية حث من الله للمؤمنين على الصبر على طاعته، وعلى احتمال ما يلقونه في امتثال أمره من المشقة ومقالة الأعداء، وعلى الاستعانة بالصلاة في الفزع إلى الله وطلب معونته.

وربط \rawi{الطبري} الآية بما قبلها من حديث القبلة وما ترتب عليها من اعتراض السفهاء، ففسر الأمر بالصبر والصلاة بأنه عون على الثبات عند ورود الأوامر الشاقة والابتلاءات التابعة لها، ثم ساق عن \rawi{أبي العالية} و\rawi{الربيع} ما يدل على أن الآية عامة في الاستعانة بالصبر والصلاة على مرضاة الله كلها.

\begin{fawaid}[معية الله للصابرين]
فسر الطبري قوله تعالى \textit{إِنَّ اللَّهَ مَعَ الصَّابِرِينَ} بمعنى: ناصرهم، ومعينهم، وظَهِيرُهم. فهذه معية خاصة تقتضي التأييد والحفظ والرعاية، وليست مجرد علم عام.
\end{fawaid}

\separator

\section{تأويل (ولا تقولوا لمن يقتل في سبيل الله أموات)}
\isnad{قال أبو جعفر:} إن الله نهى المؤمنين أن يصفوا من قُتل في سبيله بالموت الذي ينقطع معه النعيم والإدراك، وأخبر أنهم أحياء عنده في رزق وكرامة.

وساق \rawi{الطبري} آثاراً عن \rawi{مجاهد} و\rawi{قتادة} وغيرهما في أرواح الشهداء، وأنها تُرزق وتتنعم في البرزخ. ثم ناقش إشكال اختصاص الشهيد بهذا الخبر مع أن المؤمنين عموماً في نعيم البرزخ، فقرر أن الخصوصية هنا في نوع ما أُعطيه الشهيد من رزق الجنة ونعيمها قبل البعث على وجه لم يُعطه سائر المؤمنين مثله.

\begin{fawaid}[خصوصية حياة الشهداء في البرزخ]
المؤمنون عموماً منعَّمون في البرزخ، لكن الشهيد اختصه الله بفضيلة زائدة؛ وهي أنه \textbf{حي مرزوق} من نعيم الجنة على وجه أخص وأظهر من غيره، ولذلك ورد النص عليه تعظيماً لمنزلته وترغيباً في الجهاد.
\end{fawaid}

\begin{fawaid}[منهج الطبري في دفع الإشكال]
من صنيع الطبري النافع أنه لا يكتفي بحكاية المعنى المختار، بل يورد السؤال الذي قد يَرِد على ذهن القارئ ثم يجيب عنه. ففي هذا الموضع سأل: إذا كان المؤمنون جميعاً منعمين في البرزخ، فما وجه تخصيص الشهيد؟ ثم حرر الجواب من نفس الآيات والآثار.
\end{fawaid}

\separator

\section{تأويل (ولنبلونكم بشيء من الخوف والجوع...)}
\isnad{قال أبو جعفر:} هذا إخبار من الله للمؤمنين أنه مبتليهم بشيء من الشدائد: من الخوف، والجوع، ونقص الأموال، والأنفس، والثمرات، ثم أمر نبيه ﷺ أن يبشر الصابرين.

وبيّن الشيخ أن هذا الابتلاء متصل بما قبله؛ فكما ابتلاهم الله بتحويل القبلة وبالانقياد للأمر، فإنه يبتليهم كذلك بأنواع المكاره، ليتميز الصادق من المعترض.

\begin{fawaid}[الابتلاء سنة جارية بعد الإيمان]
الابتلاء ليس أمراً عارضاً في طريق المؤمن، بل هو سنة ملازمة لطريق العبودية. فكلما عظم إيمان العبد واستقامته، عُرض له من البلاء ما يمتحن به صدقه وصبره وتسليمه.
\end{fawaid}

\subsection{تأويل (الذين إذا أصابتهم مصيبة قالوا إنا لله وإنا إليه راجعون)}
\isnad{قال أبو جعفر:} هؤلاء هم الذين إذا نزلت بهم المصائب علموا أنهم عبيد لله، وأنهم راجعون إليه بعد مماتهم، فسلموا لقضائه ورضوا بحكمه.

\begin{fawaid}[معنى الاسترجاع عند المصيبة]
قول العبد: \textit{إِنَّا لِلَّهِ وَإِنَّا إِلَيْهِ رَاجِعُونَ} ليس مجرد لفظ يقال عند المصيبة، بل هو إعلان عبودية وتفويض واستسلام، يتضمن أن العبد ملك لله، وأن مرجعه إليه، فلا يعترض على ربه بل يرضى ويسلم.
\end{fawaid}

\subsection{تأويل (أولئك عليهم صلوات من ربهم ورحمة)}
فسر \rawi{الطبري} الصلوات هنا بثناء الله على عباده الصابرين وتعظيمه لهم، وضم إليها الرحمة والهداية، فصاروا هم المهتدين حقاً.

\separator

\section{تأويل آيات الصفا والمروة}
\subsection{تأويل (إن الصفا والمروة من شعائر الله)}
ذكر \rawi{الطبري} معنى الصفا والمروة، وأنهما من معالم دين الله وشعائره الظاهرة التي شرعها لعباده.

\subsection{تأويل (فلا جناح عليه أن يطوف بهما)}
استوعب \rawi{الطبري} في هذه الآية الأقوال والآثار الواردة في سبب النزول، ثم بحث حكم السعي بين الصفا والمروة، وناقش القراءة الشاذة المنقولة: \textit{فلا جناح عليه ألا يطوف بهما}، ورد الاستدلال بها من جهات متعددة:
\begin{itemize}
    \item لأنها مخالفة لرسم المصحف.
    \item ولأن سنة النبي ﷺ في المناسك بينت أن الطواف بهما مشروع معمول به.
    \item ولأن القياس العام في النسك يشهد لذلك.
\end{itemize}

\begin{fawaid}[منهج الطبري في إدارة الخلاف في آية الصفا والمروة]
من أحسن ما يلاحظ في هذا الموضع أن الطبري يبدأ بجمع الآثار، ثم يحرر موضع النزاع، ثم يذكر حجج الأقوال، ثم يرد ما لا يثبت منها بالحديث والسنة ورسم المصحف والقياس، فلا يكتفي بمجرد الحكاية.
\end{fawaid}

\subsection{تأويل (ومن تطوع خيراً فإن الله شاكر عليم)}
بيّن \rawi{الطبري} أن معنى الآية يختلف بحسب القول في حكم السعي؛ فمن قال بوجوبه حمل التطوع على ما زاد من حج أو عمرة، ومن قال بعدم الوجوب حمل التطوع على نفس السعي. ثم رجح ما يراه صواباً مع بيان وجه الآية على سائر الأقوال.

\begin{fawaid}[تفسير الآية على مقتضى كل قول قبل الترجيح]
من دقة الطبري أنه لا يقتصر على ترجيح قول بعينه، بل يبين كيف تُفهم الآية على كل مذهب معتبر في المسألة، ثم يرجح بعد ذلك. وهذا أكمل في التصوير العلمي للخلاف وأدق في فهم لوازم الأقوال.
\end{fawaid}

\separator

\section{تأويل (إن الذين يكتمون ما أنزلنا من البينات والهدى)}
\isnad{قال أبو جعفر:} نزلت هذه الآية في علماء اليهود والنصارى الذين كتموا أمر النبي ﷺ وصفته، مع أن الله بينه لهم في كتبهم.

ثم نبه الشيخ إلى أن \rawi{الطبري} مع تقريره لسبب النزول الخاص، لم يحصر الآية فيه، بل نص على أن حكمها يتناول كل من كتم علماً أوجب الله بيانه. واستشهد بما نُقل عن \rawi{أبي هريرة} رضي الله عنه في استدلاله بهذه الآية على وجوب تبليغ العلم.

\begin{fawaid}[عموم الوعيد لكل كاتم علم]
إذا نزلت الآية في صورة مخصوصة ثم دل الدليل على عموم لفظها ومعناها، لم يجز قصرها على سببها. فآية كتمان العلم نزلت في أهل الكتاب، لكنها تتناول كل من كتم علماً يجب عليه بيانه للناس.
\end{fawaid}

\separator

\section{تأويل (أولئك يلعنهم الله ويلعنهم اللاعنون)}
ذكر \rawi{الطبري} الأقوال في المراد باللاعنين؛ فقيل: هم الملائكة والمؤمنون، وقيل: البهائم والهوام وسائر الدواب. ثم رجح أن المراد بهم الملائكة والمؤمنون، مستنداً إلى ظاهر القرآن في نظائره، كقوله تعالى: \textit{أُولَٰئِكَ عَلَيْهِمْ لَعْنَةُ اللَّهِ وَالْمَلَائِكَةِ وَالنَّاسِ أَجْمَعِينَ}.

وأجاز من حيث الاحتمال أن تلعنهم البهائم وسائر الخلق، لكنه لم يجعل ذلك هو تفسير الآية لعدم الحجة القاطعة عليه.

\begin{fawaid}[الفرق بين الاحتمال والتفسير]
ليس كل معنى \textit{محتمل} يجوز جعله هو \textit{تفسير} الآية. فقد يحتمل العموم دخول أفراد كثيرة تحته، لكن تعيين فرد مخصوص على أنه المراد يحتاج إلى حجة. ولهذا لم يخصص الطبري ``اللاعنين'' بالهوام والبهائم مع إمكان دخولها في العموم.
\end{fawaid}

\begin{fawaid}[أمور الغيب لا يقطع فيها بمجرد الإمكان]
ما كان من أمور الغيب لا يجوز الجزم بتعيينه لمجرد وروده احتمالاً أو لمناسبة عقلية، بل لابد فيه من خبر تقوم به الحجة. ولذلك قبل الطبري إمكان دخول بعض الخلق في اللعن، لكنه امتنع من القطع بأنهم المرادون بالآية من غير دليل قاطع.
\end{fawaid}

\separator

\section{تأويل (إلا الذين تابوا وأصلحوا وبينوا)}
بيّن \rawi{الطبري} أن التوبة هنا ليست مجرد رجوع قلبي، بل لابد معها من إصلاح الحال وبيان ما كان العبد كتمه من الحق؛ لأن الذنب تعلق بإخفاء البيان، فلا تتم التوبة منه إلا بإظهاره.

\begin{fawaid}[التوبة من كتمان العلم لا تتم إلا بالبيان]
من كتم حقاً أو لبس على الناس في دين الله لا يكفيه الندم المجرد، بل من تمام توبته أن يصلح ما أفسده وأن يبين ما كان قد أخفاه؛ لأن الذنب تعلق بالناس كما تعلق بحق الله.
\end{fawaid}

\separator

\section{تأويل (إن الذين كفروا وماتوا وهم كفار...)}
ذكر \rawi{الطبري} أن هؤلاء عليهم لعنة الله والملائكة والناس أجمعين، وأنهم خالدون في هذا العذاب لا يخفف عنهم ولا يُنظرون، وفيه رد على من زعم انقطاع عذاب الكفار أو خروجه عن مقتضى الخلود.

\separator

\section{تأويل (وإلهكم إله واحد لا إله إلا هو الرحمن الرحيم)}
فسر \rawi{الطبري} هذه الآية بإثبات توحيد الله سبحانه، وأنه المتفرد بالإلهية المستحق للعبادة دون ما سواه، ثم أتبعها بآيات الكون الدالة على ربوبيته ووحدانيته.

\subsection{تأويل (إن في خلق السماوات والأرض...)}
ساق \rawi{الطبري} دلائل الآفاق والأنفس في خلق السماوات والأرض، واختلاف الليل والنهار، والفلك، والمطر، وبث الدواب، وتصريف الرياح، والسحاب المسخر، مبيناً أنها آيات لقوم يعقلون.

\separator

\section{تأويل (ومن الناس من يتخذ من دون الله أنداداً)}
ذكر \rawi{الطبري} أن هؤلاء جعلوا لله أنداداً يحبونهم كحب الله، وأن المؤمنين أشد حباً لله. ثم شرح ما يقع يوم القيامة من انكشاف حقيقة الضعف والخذلان لمن اتبع غير الله.

\subsection{تأويل (إذ تبرأ الذين اتبعوا من الذين اتبعوا)}
اختلفوا في المتبوعين: أهم الرؤساء والكبراء، أم الشياطين؟ فرجح \rawi{الطبري} العموم، وأن كل متبوع على الكفر والضلال يتبرأ يوم القيامة من أتباعه، مع بقاء دلالة السياق على أن الأنداد المذكورين يدخل فيهم رؤوس الضلال من الرجال.

\subsection{تأويل (وتقطعت بهم الأسباب)}
جمع \rawi{الطبري} الأقوال في معنى الأسباب: من المودة، والمنازل، والأرحام، والأعمال، وسائر الوصل. ثم رجح أن الآية تعم جميع الأسباب التي كان الكفار يتسببون بها في الدنيا إلى مقاصدهم، فتتنقطع كلها يوم القيامة.

\begin{fawaid}[حمل العموم على جميع معانيه إذا لم توجد حجة التخصيص]
إذا دل اللفظ على معنى عام، ووردت تحته أفراد متعددة لا تعارض بينها، ولم تقم حجة على التخصيص، فالأولى إبقاؤه على عمومه. ولذلك رجح الطبري أن ``الأسباب'' تشمل كل ما كانوا يتوصلون به في الدنيا من مودة وقرابة وعمل واتباع.
\end{fawaid}

\subsection{تأويل (كذلك يريهم الله أعمالهم حسرات عليهم)}
ذكر \rawi{الطبري} قولين مشهورين:
\begin{itemize}
    \item أن المراد: الأعمال التي كانوا قد عملوها فصارت عليهم ندامة وحسرة.
    \item أن المراد: الأعمال الصالحة التي كان ينبغي أن يعملوها فحُرموها.
\end{itemize}

ورجح الأول؛ لأنه الموافق لظاهر اللفظ، أما الثاني فمع كونه محتملاً فهو بعيد ولا تنهض به حجة تترك لأجلها دلالة الظاهر.

\begin{fawaid}[الظاهر مقدم على الباطن المحتمل]
من قواعد الطبري أن ظاهر الآية لا يُعدل عنه إلى معنى باطن بعيد إلا بحجة قوية. فإذا احتمل اللفظ وجهاً خفياً لكنه خالف الظاهر ولم يقم عليه دليل، كان الواجب إبقاء الآية على ظاهرها.
\end{fawaid}

\separator

\section{تأويل (يا أيها الناس كلوا مما في الأرض حلالاً طيباً)}
أمر الله الناس أن يأكلوا من الحلال الطيب الذي أحله لهم، ونهاهم عن اتباع خطوات الشيطان فيما يزينه من تحريم ما لم يحرمه الله أو تحليل ما حرمه الله.

وجمع \rawi{الطبري} في تفسير ``خطوات الشيطان'' أقوال السلف: بأنها أعماله، وخطاياه، وآثاره، وطاعته، والنذور في المعاصي. ثم قرر أنها متقاربة ترجع إلى اتباع آثاره ومسالكه.

\begin{fawaid}[خطوات الشيطان اسم جامع لآثاره ومسالكه]
ليس المراد بخطوات الشيطان فعلاً واحداً بعينه، بل كل طريق يقود إلى معصية الله أو تبديل شرعه أو طاعة الشيطان في أثر من آثاره؛ فكل ذلك داخل في النهي.
\end{fawaid}

\subsection{تأويل (إنما يأمركم بالسوء والفحشاء وأن تقولوا على الله ما لا تعلمون)}
بيّن \rawi{الطبري} أن الشيطان لا يأمر إلا بما فيه السوء والفحشاء والقول على الله بغير علم، فجعل أعظم أبوابه الافتراء على الله في التحليل والتحريم والدين.

\separator

\section{تأويل (وإذا قيل لهم اتبعوا ما أنزل الله)}
ذكر \rawi{الطبري} حال من عارض الوحي بتقليد الآباء، وقرر أن من لا عقل له هو الذي يتبع الجاهل لمجرد كونه أباً أو متبوعاً، دون أن ينظر في الهدى والحق الذي جاء من عند الله.

\separator

\section{تأويل (ومثل الذين كفروا كمثل الذي ينعق...)}
استعرض \rawi{الطبري} قولين مشهورين في الآية:
\begin{itemize}
    \item أن مثل الكافرين مع واعظهم كمثل البهائم التي تسمع صوت الناعق ولا تفقه كلامه.
    \item أن مثلهم في دعاء آلهتهم كمثل من ينعق بما لا يسمع إلا صوتاً مجرداً.
\end{itemize}

ورجح القول الأول، لأنه الأليق بالسياق؛ إذ الآيات في اليهود الذين يكتمون الحق ويعرضون عن الموعظة، لا في عبدة الأصنام.

\begin{fawaid}[السياق من أقوى قرائن الترجيح]
من أعظم ما يعتمد عليه الطبري في التفسير النظر إلى السياق السابق واللاحق. فلما كان السياق هنا في أهل الكتاب المعرضين عن الحق، رجح أن المثل في ضعف فهمهم للموعظة، لا في دعائهم الأصنام.
\end{fawaid}

\subsection{تأويل (صم بكم عمي فهم لا يعقلون)}
فسر \rawi{الطبري} هذه الأوصاف بأنها أوصاف معنوية: صم عن سماع الحق سماع قبول، بكم عن النطق به وبيانه، عمي عن إبصاره واتباعه. فليست الآفة في الجوارح، وإنما في عدم الانتفاع بها.

\separator

وقف الشيخ هنا، بعد أن أتم في هذا اللقاء جملة كبيرة من آيات الصبر والابتلاء، وآيات الشعائر، وكتمان العلم، والتوحيد، وذم التقليد، والأمثال المضروبة للكافرين. وختم بتكليف الطلاب بتلخيص منهج \rawi{الطبري} في آية الصفا والمروة، وتحرير الأقوال في مثل الذين كفروا كمثل الذي ينعق بما لا يسمع إلا دعاء ونداء.
